\section{Related Work}

\paragraph{Neural scaling laws}
The study of scaling laws for language models was popularized by \citet{kaplan2020scaling}, who showed that loss decreases as a power law in model size, dataset size, and compute. \citet{hoffmann2022training} refined these findings with the Chinchilla scaling laws, establishing compute-optimal training recipes that balance model size and data. Relating to our work, \citet{muennighoff2023scaling} generalize Chinchilla to a single data-constrained domain, finding that repeating data up to approximately four epochs causes negligible degradation, with meaningful gains extending to roughly 16 epochs and ending at 40. They propose scaling laws that model the diminishing value of repeated tokens. Subsequent work has extended scaling laws to downstream task performance~\citep{wei2022emergent}, mixture-of-experts architectures~\citep{abnar2025parameters,krajewski2024scaling,wang2024scaling}, continued pretraining~\citep{que2024d,liew2025acceleration,seto2026optimal}, finetuning~\citep{bethune2025scaling,zhang2024scaling}, and data quality~\citep{chang2024scaling}.
Our work extends this line of research providing a predictive scaling law for data constrained mixtures, an under-explored regime where only one component of a data mixture is unlimited while others are severely constrained.

\paragraph{Data mixing for language models}
Training on a mixture of data domains has become standard, with corpora aggregated from multiple web sources~\citep{gao2020pile,cerebras2023slimpajama,soldaini2024dolma} with varying data size per domain, and recent models are trained on carefully tuned domain mixtures~\citep{bakouch2025smollm3,olmo2025olmo}. Prior work also focus on optimizing domain weights using distributionally robust optimization  \citep{xie2024doremi,fan2024doge}, or by discovering domains through clustering the pretraining corpus~\citep{diaonemotron,grangier2025task}.
The domain mixture scaling law (DMSL) framework~\citep{ye2024data,shukor2025scalinglawsoptimaldata} models how per-domain loss depends on mixture weights and total compute. Our repetition-aware scaling law builds on DMSL by incorporating the effect of data repetition, which becomes critical when any mixture component is data-limited. Finally, \citet{anonymous2026mixdonttune} compare data mixing against hyperparameter tuning (weight decay, learning rate) as competing strategies for data-constrained pretraining.
The focus of our work is on how to set the mixture optimally once mixing is chosen by characterizing how the mixing ratio, repetition, and model size jointly determine target-domain loss, and that eliminates the need for per-configuration sweeps.


\paragraph{Data-constrained pretraining}
Several works study data-constrained pretraining. \citet{hernandez2022scaling} propose parametric models for the value of repeated data in multi-epoch training, and \citet{goyal2024scaling} compare repeating high-quality data against training on fresh lower-quality data, concluding that the benefit of filtering depends on the total compute budget. When high-quality data for the target domain is constrained, synthetic data has been used as an alternative to repetition. \citet{gunasekar2023textbooks} pretrained on synthetic ``textbook quality'' data, and \citet{maini2024rephrasing} demonstrate that rephrasing web data offer alternatives to repeating data. For multi-domain settings, \citet{seto2025training} study pretraining bilingual models when target-language data is constrained, showing that high-quality auxiliary-language data can partially substitute for target-language data in typologically close languages, but that model scaling has diminishing returns. Low-resource language modeling faces similar constraints, where mixing with high-resource languages is a common strategy ~\citep{joshi2020state}. This work studies different scenarios where data is constrained in at least one target domain, and unconstrained in other generic data, and provides a recipe for determining the target domain loss.  This is tangential to other approaches which aim to improve the diversity of data by building more target data, or incorporating more generic data.





